\section*{Ethics Considerations}\label{sec:ethics}
Our SoK paper reviewed and systematized existing web tracker detection studies since the 1990s. While related work may not have discussed significant ethical issues, we acknowledge that it is insufficient to apply the analyses from past publications to our research. Therefore, we conduct an ethical evaluation in light of current technologies, societal expectations, and data protection requirements using the framework provided by the Menlo Report on ``Ethical Principle Guiding Information and Communication Technology Research''~\cite{menlo_report_2012}. 

\noindent\textbf{Stakeholders} We identify the following primary stakeholders of our research: 

\begin{enumerate}[itemsep=0.3mm, parsep=0pt]
    \item \textbf{Authors} of the reviewed studies, whose work or artifacts were analyzed and whose reputation or research trajectory could be influenced by our interpretation. 
    \item \textbf{Individuals} whose personal data may be contained in datasets used in the reviewed studies, even if anonymized, as their privacy could be indirectly affected by how such data is represented. 
    \item The \textbf{academic research community}, who may rely on our systematization of the current literature for future work. 
    \item The \textbf{general public}, given the broader societal implications of web tracking and privacy research. 
\end{enumerate}

\noindent\textbf{Potential Harms.} Considering the interests of all primary stakeholders, we recognize several potential harms, both tangible and intangible, that may occur in the short and long term as a result of our paper and its publication. Firstly, our research and experiments were conducted to the benefit of our community and do not pose any harm to these authors, nor with any intention to discredit their work (\textit{Respect for Persons}). However, there remains a possibility of reputational damage if our analysis and  interpretations are misrepresented or taken out-of-context. Secondly, although we treated all identified primary and supplementary studies impartially without bias toward any population, group, or person (\textit{Justice}), our eligibility criteria did not include an assessment of the ethical conduct of each study. Thus, it is possible that some included work did not meet current ethical standards. This risk is compounded by the fact that some studies are older and may have been conducted under ethical norms and data protection regulations that differ from current norms. Therefore, it is necessary to recognize the possibility that our paper indirectly legitimizes prior unethical research if it took place. While we trust our colleagues adhere to ethical research practices (and to the best of our knowledge, we have no reason to believe otherwise), this paper, nonetheless carries with the obligation to critically consider the ethical integrity of the studies it engages with. Finally, our systematization of detection approaches could be misused by adversaries to evade detection. 

\noindent\textbf{Mitigations and Unmitigated Risks.} To minimize these risks, we confined our systematization to publicly available sources. We did not recruit or interact with human participants, nor did we gather identifiable information about individuals. Experiments were only ran on code and data that was publicly released by the authors for others to use and we cited all sources. We did not disclose any implementation details beyond what authors already released. We also plan to inform authors post hoc and remain available to clarify our interpretations to reduce potential misrepresentations. Notwithstanding these measures, residual risks remain, such as the absence of an explicit ethics-screening criterion (and analysis), the inability to verify the compliance of older studies with current ethical and legal standards, and the potential for adversarial adaptation of tracking techniques. We acknowledge these risks and intent to address them in future work by developing explicit ethics assessment criteria and guidelines. 

\noindent\textbf{Ethical Analysis} From a consequentialist perspective, we have to weigh the anticipated benefits and harms against one another. The principal benefit is a holistic understanding of detection, which can guide future research and highlight new directions for advancing privacy-enhancing technologies. We argue that these contributions to the academic community, and the broader societal value of systematized privacy knowledge, outweigh the probable risks identified above. We nevertheless acknowledge that the full range of potential consequences cannot be entirely foreseen. 

From a deontological perspective, we must assess whether our actions adhered to the principles of \textit{Beneficence}, \textit{Respect for Persons}, \textit{Justice}, and \textit{Respect for Law and Public Interest}. Our study did not involve human subjects and posed no direct violations of rights, moreover, there is neither disclosure of new vulnerabilities or user data, nor probing of live web services. However, the absence of an explicit ethics-screening step in our methodology means that our moral duty to respect the rights and interests of individuals represented in underlying studies remains only partially fulfilled. We see this as an area of improvement for our future research. At the same time, we upheld other duties by relying exclusively on publicly available artifacts, attributing all sources, avoiding the disclosure of new vulnerabilities, and striving for impartiality in our selection and treatment of studies. Finally, to the best of our knowledge, we complied with regulations, copyrights, and applicable laws. 

\noindent\textbf{Decision to Proceed}
After considering both consequentialist and deontological perspectives, we conclude that conducting and publishing this study constitutes an ethically justified decision. We expect the benefits of this work to outweigh the potential risks and to serve not only the academic community but also the original authors of the reviewed studies and the public at large. Our primary motivation has been to advance knowledge in the domain of web tracker detection and web privacy, and we remain committed to pursuing this goal in the most responsible and transparent manner possible. To foster trust, independent verification, and practical reuse, we release all of our artifacts openly. Finally, we acknowledge the residual risks identified in this analysis and intend to address them in future work by incorporating explicit ethics eligibility criteria into the process of literature inclusion.
